\chapter{使用データと解析条件}\label{ch:data}
\section{使用データの出典と本研究の対象}
本研究は，植木研究室内の他の研究メンバーが実施したプール沸騰実験の既存音響・熱計測データを用いる解析研究である．主比較には，保存資料上で2025年6月11日および6月18日に対応づけられたデータを使用する．著者はこれらの実験・収録を実施しておらず，既存データを用いる解析を担当する．個々の実験実施者の氏名は現時点で特定していないが，データの由来と著者の担当を以上のように区別する．
\need{D02}{データ取得，熱流束算出，音響前処理，モデル実装・実行，評価・考察について，著者・共同研究者・既存コードの担当を確認する．本稿の解析上の貢献と，著者自身が実施した作業の範囲は最終稿で対応づける．}

本稿で扱う解析上の内容は，学習と評価の条件整理，異なる表現を用いる回帰器の比較，統合による誤りの変化の分析，入力依存の検討である．既存実験の方法は出典を付して説明し，文献で確認した条件と対象日の記録から確認した条件を区別する．

\section{音響・熱計測データの概要}
\subsection{対象条件と収録単位}
データはサブクール度20\,$^\circ$C，細線径0.3\,mmに対応する保存系列に置かれている．解析仕様\cite{localData}による音響サンプリング周波数は44,100\,Hzであり，1録音の60秒を1秒ずつ60区間へ分割する．本稿では元の録音ファイルをWAV，1秒区間をchunkと呼ぶ．chunkは音響入力の単位であり，独立な沸騰実験の回数ではない．

\begin{table}[htbp]
\centering\small
\caption{主比較に用いるデータの構成．熱流束は保存ラベルの範囲である．}
\label{tab:data}
\begin{tabular}{lrrr}
\toprule
収録日 & WAV数 & 1秒区間数 & 熱流束範囲 [\kw{}] \\
\midrule
2025年6月11日 & 18 & 1,080 & 約0--859.34 \\
2025年6月18日 & 18 & 1,080 & 約0--873.64 \\
合計 & 36 & 2,160 & 約0--873.64 \\
\bottomrule
\end{tabular}
\source{保存予測の正解ラベルと主実行記録\cite{localMainRun}．}
\end{table}

表\ref{tab:data}の範囲は，現データに含まれる記録範囲を示す．最大記録熱流束を物理的な限界熱流束（CHF）と同一視しない．また，無加熱付近のラベルには計算丸めに相当するごく小さい値が含まれるため，範囲表では約0と表記した．ラベルを音響の大きさから作り直したり，静かな区間を未沸騰へ変更したりはしていない．

\subsection{原実験条件と本稿で確認できた範囲}
統合データの\path{source_metadata/}には，各日の電気計測CSV，表計算ファイルおよびロガー出力の計6ファイルが保存されている．元の日別フォルダの同名6ファイルとSHA-256が全て一致することを確認した\cite{localMetadata}．CSVヘッダーには所有者欄\texttt{UekiLab}，ロガー型式34972A，水温・シャント電圧・電極電圧等の計測項目があり，対象日の電気計測資料へ遡れる．ただし，所有者欄は実験実施者の氏名を示すものとは限らず，ロガー型式も音響収録器の型式ではない．この確認だけでは，録音区間と熱流束算出区間の対応，装置条件およびONBの観察根拠までは確定しない．

柴崎\cite{shibasaki2025}の第2章では，白金細線径0.3\,mm，長さ40\,mm，ハイドロフォンによる44,100\,Hzの音響収録，印加電圧段階ごとの60秒録音を記述している．手塚\cite{tezuka2026}も白金細線を用いた水のサブクールプール沸騰装置と音響収録を記述している．これらは同研究室の既往実験手法の根拠であり，各条件が2025年6月の主対象日にそのまま適用されたかは，対象日の記録との対応を区別して扱う．

\need{D03}{対象日の細線長・支持条件，作動流体と水量，圧力・水温，センサ型番・位置，収録ゲイン，校正の有無を原資料で確定する．同じ0.3\,mm系列であることから全条件の一致を推測しない．}

\section{熱流束とONBラベルの根拠}
\subsection{熱流束ラベルの意味}
各WAVには一つの熱流束段階が対応し，その60個のchunkへ同じラベルを割り当てる．これは，録音段階に対応する熱流束を回帰の正解とする設定である．録音中の一秒ごとの熱流束を別々に測定したという意味ではなく，その一秒に気泡が発生しているかというイベント正解とも異なる．同じ段階で音響が弱いchunkがあっても，熱流束の存在自体が否定されるわけではない．

既往実験手法の3.5.1節\cite{shibasaki2025}には，白金細線に加わる電圧$V_{\mathrm{Pt}}$と電流$I$から，
\begin{equation}\label{eq:heat_flux_source}
 q=\frac{V_{\mathrm{Pt}}I}{\pi dL},\qquad
 I=\frac{V_{\mathrm{shunt}}}{R_{\mathrm{shunt}}}
\end{equation}
により熱流束を求める方法が記されている．同文献の値は，細線径$d=0.3$\,mm，長さ$L=40$\,mm，シャント抵抗$R_{\mathrm{shunt}}=1.667\times10^{-3}$\,$\Omega$である．$V_{\mathrm{shunt}}$はシャント電圧であり，$\pi dL$は細線側面の伝熱面積を表す．これらは記載文献の実験条件であり，対象2日への適用値を未照合のまま確定したものではない．

同文献の付録A.1には，60秒の録音時にロガーで記録したシャント電圧と白金電圧のうち，定常状態となった範囲を平均して熱流束等を算出したと記されている．従って，既往手法の時間平均の考え方は確認できる．対象2日の具体的な平均区間と各WAVの対応は，この方法の記載とは別に扱う．
\need{D04}{算出式と定常平均の方法は文献で確認済みである．残る確認は，対象2日のラベル作成時に用いた抵抗値・有効寸法，ロガーの平均行範囲，CSVと録音の対応である．不確かさを定量的に示す場合は，校正値・寸法誤差・計測誤差の根拠も必要となる．}

\subsection{ONB閾値の値と出典の段階}
本研究では，目視で核沸騰の蒸気泡発生を初めて確認した際の熱流束をONBとする．この操作的定義は，既往実験手法の4.1.2節\cite{shibasaki2025}に記載された，各電圧段階の1分間動画を分析して最初の蒸気泡を確認し，対応する電圧からONB熱流束を求める方法と一致する．主比較には日別のONB熱流束を用い，現在使用している値を表\ref{tab:onbthreshold}に示す．これらは2026年9月24日の本人確認を記録した現行値\cite{localONB}である．判定の定義と採用値は以上のように説明できるが，対象2日の個別観察資料と音響内の発生秒の対応までは確認していない．
\begin{table}[htbp]
\centering
\caption{日別ONB閾値として現在採用している値．}
\label{tab:onbthreshold}
\begin{tabular}{lrr}
\toprule
収録日 & 保存値 [\wm{}] & 表示値 [\kw{}] \\
\midrule
2025年6月11日 & 221,505.1102 & 221.5051102 \\
2025年6月18日 & 271,677.6816 & 271.6776816 \\
\bottomrule
\end{tabular}
\end{table}

過去の自動解析には，抵抗と熱流束の関係において直線性が失われる候補点をONBと呼んだ出力があったが，その値は現在のONBとは異なる．本稿ではその自動候補を正解へ戻さない．また，2日分のONBの算術平均246.5913959\,\kw{}を用いる旧表示があるが，主比較の見逃し・誤報および到達熱流束には，平均閾値ではなく表\ref{tab:onbthreshold}の日別閾値を適用する．

目視でONB熱流束を定めることと，音響の何秒目に最初の気泡が発生したかを記録することは異なる．後者には動画等と音響の時刻対応が必要である．本稿は熱流束段階に対する判定を評価し，発生時刻からの検知遅延を秒で評価したものではない．
\need{D05}{目視によるONBの定義と採用値は確認済みである．対象日の観察記録・対応段階が残っていれば出典を補う．判定の分解能や不確かさを数値で述べるには追加の根拠が必要である．同期した発生秒の記録の有無は未確認だが，現在の段階判定をまとめるための必須取得事項にはしない．}

\section{解析対象データの構成}
\subsection{学習区間と評価区間}
主比較では，各WAVの60区間をshuffleによって45学習区間と15評価区間へ分ける．分割seedは42であり，学習は$36\times45=1,620$区間，評価は$36\times15=540$区間となる．同じ評価区間IDを全モデル・全雑音条件で共有する．ピーク高さなどによる区間選別は行わない．

この分割では，同じWAVに属する兄弟chunkが学習側と評価側の両方に含まれる．従って，評価の対象は既知の収録条件・熱流束段階に含まれる未使用区間であり，未知WAVや未知日への外挿を試したものではない．区間の重複を防いでいることと，録音単位の独立性を確保していることは異なる．本稿はこの条件を主比較の成立範囲として明記する．

主評価集合には，日別ONB未満が225区間，ONB以上が315区間ある．日別ONBの$\pm10\%$を近傍とすると30区間となり，現在の離散的な熱流束段階では，両日の記録上ONB段階にある各15区間だけが該当する．つまり，この近傍評価にはONB未満側の近接段階は含まれない．「ONBの両側を均等に評価した帯域」とは呼ばない．

\subsection{雑音条件と対象範囲}
無雑音入力に加え，水流音を付加した基準SNR $0,-4,-8,-12,-16,-20$\,dBの6条件を用いる．ここで無雑音とは人工的な水流音を付加しない条件であり，元録音に背景音が全くないことを意味しない．7条件は同じ録音を加工した対応のある比較であり，独立な7回の実験ではない．基準SNRと区間ごとの実現SNRの違いは第\ref{ch:methods}章で定義する．

主結果には10月15日の完了実行を用い，10月14日の誤り分解，10月7日の入力対照を補助証拠として参照する．途中停止した条件別雑音学習の実行や，2026年の新しい録音日の未評価結果を本稿へ自動的に加えない．一方，既存の異条件の診断から評価範囲を考察する場合は，主系列の数値とは区別して付録に示す．

本章により，研究が答えられるのは，日別の記録上ONBを基準とした既知WAV内の回帰・判定であることを定めた．次章では，このデータと正解の範囲で，どの入力・モデル・統合・指標を用いるかを定義する．
